\textbf{Commun à tous les candidats}


















\begin{enumerate}
\item 








Pour $t = - 2$, on trouve les coordonnées de B.
\item 
Le vecteur $\vect{\text{AB}}$ admet pour coordonnées : $\begin{pmatrix}2 - 1~;~1 - 0~;~0 - 2\end{pmatrix}$ soit $\begin{pmatrix}1~;~1~;~- 2\end{pmatrix}$








\item 









$M(x~;~y~;~z) \in (\text{AB})$ s'il existe $t \in \R$ tel que $\vect{\text{A}M} = t\vect{\text{AB}}$ soit :

$\left\{\begin{array}{l !{=} l}
x - 1& 1\times t\\
y - 0& 1\times t\\
z - 2& - 2\times t
\end{array}\right.\, t \in \R \iff 
\left\{\begin{array}{l !{=} r}
x & 1 +  t\\
y &t\\
z & 2 - 2t
\end{array}\right.\, t \in \R 
$.

En posant $t = 1 - u$, on obtient :

$M(x~;~y~;~z) \in (\text{AB}) \iff \left\{\begin{array}{l !{=} r}
x & 2 - u\\
y &1 - u\\
z & 2u
\end{array}\right.\, u \in \R$. Donc réponse B.
\item 







La droite $\Delta$ a pour vecteur directeur $\delta\begin{pmatrix}2\\1\\-1\end{pmatrix}$

Soit $\mathcal{P}$ le plan dont on cherche une équation.

$M(x~;~y~;~z) \in \mathcal{P} \iff \vect{\text{C}M} \cdot \delta = 0 \iff 2(x - 0) + 1(y - 1) -1(z - 2) = 0 \iff$

$ 2x + y - 1 - z + 2 = 0 \iff 2x + y - z + 1 = 0$.
\item 







En faisant apparaître le point A dans chaque vecteur (Chasles), on obtient :

$\vect{\text{OA}} + \vect{\text{AD}}  = 3\vect{\text{OA}}  - \vect{\text{OA}}  - \vect{\text{AB}}  - \vect{\text{OA}}  - \vect{\text{AC}} \iff \vect{\text{AD}}  = - \vect{\text{AB}} - \vect{\text{AC}} $ : le vecteur $\vect{\text{AD}}$ est une combinaison linéaire des vecteurs $\vect{\text{AB}}$ et $\vect{\text{AC}}$, ces trois vecteurs sont donc coplanaires.
\end{enumerate}

\bigskip
